\section*{第 \ref{ch:b1:observing-the-sky} 章——望远镜、探测器与测光}

\begin{solution}{exo:b1:observing-the-sky:1}
星等相差 5 等，即亮度相差 100 倍：星等为 1.0 的恒星
发出的通量多 100 倍。参宿七与北极星：$10^{0.4\times(2.02 - 0.13)} =
10^{0.756} = 5.70$；参宿七星等较小，因而更亮。
\end{solution}

\begin{solution}{exo:b1:observing-the-sky:2}
$d = 1/0.19495 = \qty{5.13}{\parsec}$，即 $5.1295\times206265 =
\qty{1.058e6}{\au}$，以及 $5.1295\times3.2616 = \qty{16.7}{\lightyear}$。
\end{solution}

\begin{solution}{exo:b1:observing-the-sky:3}
$1.22\times\num{550e-9}/2.4 = \qty{2.80e-7}{rad} = \ang{;;0.058}$，
约为\omterm{def:b1:observing-the-sky:seeing}{视宁度} \ang{;;0.40} 的七分之一：在太空中，望远镜
达到衍射极限。若口径为 \qty{6.5}{m}、波长为 \qty{2.0}{\micro\metre}：
$1.22\times\num{2.0e-6}/6.5 = \qty{3.75e-7}{rad} = \ang{;;0.077}$.
\end{solution}

\begin{solution}{exo:b1:observing-the-sky:4}
$B - V$ 自参宿七（$-0.03$）升至牵牛星（0.22）、北极星（0.60）及
参宿四（1.85）：最热者为参宿七，其次为牵牛星、北极星和
参宿四。距离使两个波段变暗的倍数相同，因此 $B - V$
不受其影响（只要途中没有尘埃使光变红）。
\end{solution}

\begin{solution}{exo:b1:observing-the-sky:5}
$X_1 = \sec\ang{30} = 1.1547$, $X_2 = \sec\ang{60} = 2.0000$.
$k = (6.590 - 6.412)/(2.0000 - 1.1547) = 0.178/0.8453 = \num{0.211}$ 每
\omterm{def:b1:observing-the-sky:extinction}{大气质量}，且 $m_0 = 6.412 - 0.2106\times1.1547 = 6.169$。
\end{solution}

\begin{solution}{exo:b1:observing-the-sky:6}
$\mathrm{SNR} = \sqrt{400} = 20$。SNR 随 $\sqrt t$ 增长：若取 100，
$t = \qty{10}{s}\times(100/20)^2 = \qty{250}{s}$。误差为 \num{0.01}
星等，则需 $\mathrm{SNR} = 1.086/0.01 = 108.6$，因此
$t = 10\times(108.6/20)^2 = \qty{295}{s}$.
\end{solution}

\begin{solution}{exo:b1:observing-the-sky:7}
距离为 $d = 1/0.00378 = \qty{264.6}{\parsec}$，\omterm{def:b1:observing-the-sky:absolute}{绝对
星等} $M = 0.13 - 5\log_{10}26.46 = 0.13 - 7.11 = -6.98$。$L/L_\odot = 10^{0.4\times(4.83 + 6.98)} =
10^{4.724} = \num{5.3e4}$（V 波段）。
\end{solution}

\begin{solution}{exo:b1:observing-the-sky:8}
$\pi\times1.2^2 = \qty{4.52}{m^2}$，而 $25/4.52 = 5.5$：韦伯望远镜
在短 5.5 倍的曝光时间内收集到相同数目的光子。若整个
直径为 \qty{6.5}{m} 的圆盘，面积为 $\pi\times3.25^2 = \qty{33.2}{m^2}$；
实际镜面仅占其 $25/33.2 = \qty{75}{\percent}$，原因在于 18 块六边形子镜
未覆盖该圆盘的边缘，子镜间的缝隙与
中心开口均不收集光线。
\end{solution}

\begin{solution}{exo:b1:observing-the-sky:9}
(a) $\qty{10}{\parsec} = \qty{3.0857e17}{m}$，且 $F = \num{3.828e26}/(4\pi
\times(\num{3.0857e17})^2) = \num{3.828e26}/\num{1.1965e36} =
\qty{3.20e-10}{W.m^{-2}}$，$1361/\num{3.20e-10} = \num{4.3e12}$ 倍，
比地球处的通量少这么多倍。
(b) $5\log_{10}(d/\qty{10}{\parsec}) = -26.74 - 4.83 = -31.57$，因此
$d = \qty{10}{\parsec}\times10^{-6.314} = \qty{4.853e-6}{\parsec}$，即
$\num{4.853e-6}\times206265 = \qty{1.001}{\au}$.
\end{solution}

\begin{solution}{exo:b1:observing-the-sky:10}
大气层外的通量为 $F_0$；在地面则为 $F = F_0e^{-\tau_0X}$。于是
$m - m_0 = -2.5\log_{10}(F/F_0) = 2.5\log_{10}e\;\tau_0X$，即 $kX$，其中
$k = 2.5\log_{10}e\;\tau_0 = 1.0857\,\tau_0$。当 $k = 0.160$ 时：
$\tau_0 = 0.160/1.0857 = 0.147$。\omterm{def:b1:spherical-astronomy:celestial-sphere}{天顶}处，到达的光占 $e^{-0.147} = 0.863$
，即 \qty{86}{\percent}；在 $X = 2.31$ 处，$e^{-0.147\times2.31} =
e^{-0.340} = 0.712$，即 \qty{71}{\percent}。当 $z = \ang{85}$ 时，$\sec z =
11.5$；但平行平面模型使地平线处的光程为无穷大，
而弯曲大气中的光程保持有限，且折射
也使光线弯折：在地平线附近，该公式高估了
\omterm{def:b1:observing-the-sky:extinction}{大气质量}，且其误差增长很快。
\end{solution}

\begin{solution}{exo:b1:observing-the-sky:11}
(a) $\mathrm{SNR} = Rt/\sqrt{(R + B)t}$，故 $t = \mathrm{SNR}^2(R + B)/R^2
= 100\times2050/2500 = \qty{82}{s}$。(b) \omterm{def:b1:observing-the-sky:seeing}{视宁度}减半，使像的
面积（从而 $B$）除以 4：$B = 500$，且 $t = 100\times550/2500 =
\qty{22}{s}$。(c) 当 $B \gg R$ 时，$t \approx \mathrm{SNR}^2B/R^2$。源的
计数率 $R$ 随 $D^2$ 增长；天空计数率 $B$ 随集光面积
$D^2$ 与像的天空面积 $\theta^2$ 之积增长。因此 $t \propto
\theta^2D^2/D^4 = \theta^2/D^2$：\omterm{def:b1:observing-the-sky:seeing}{视宁度}减半的作用，
相当于口径加倍。
\end{solution}

\begin{solution}{exo:b1:observing-the-sky:12}
(a) $N = 1200\times100 = \num{120000}$ 条；一级光谱中
分辨本领为 \num{1.2e5}，且 $\Delta\lambda = 550/(\num{1.2e5}) =
\qty{4.6e-3}{nm}$。(b) $v = c\,\Delta\lambda/\lambda = c/(\num{1.2e5}) =
\qty{2.5}{km.s^{-1}}$。(c) $\Delta\lambda = 550\times29.78/\num{299792} =
\qty{0.055}{nm}$，约相当于 12 个分辨单元。观测者须计算
对观测的日期与时刻，求出地球公转速度朝向恒星的分量，
并减去它，以得到恒星
相对于太阳的速度。
\end{solution}

\begin{solution}{pb:b1:observing-the-sky:1}
\textbf{1.} $d = 1/0.37921 = \qty{2.637}{\parsec}$.

\textbf{2.} $2.637\times206265 = \qty{5.44e5}{\au}$;
$2.637\times\num{3.0857e16} = \qty{8.14e16}{m}$.

\textbf{3.} $2.637\times3.2616 = \qty{8.60}{\lightyear}$：光离开
天狼星已是 8.6 年前。

\textbf{4.} 方向摆动 $2p = \ang{;;0.758}$，约为
\omterm{def:b1:observing-the-sky:seeing}{视宁度}的两倍。摆动量是从像中心的位置测得，
而非自边缘算起：对大量光子取平均后，中心可被确定到
像宽度的很小一部分，计数的光子越多，
定位就越精确。

\textbf{5.} $\Delta d/d = \Delta p/p = 0.001/0.37921 = \qty{0.26}{\percent}$,
即 $8.60\times0.0026 = \qty{0.023}{\lightyear}$。

\textbf{6.} $10^{0.4\times(0.03 + 1.46)} = 10^{0.596} = 3.94$.

\textbf{7.} $10^{0.4\times(26.74 - 1.46)} = 10^{10.112} = \num{1.29e10}$.

\textbf{8.} $10^{0.4\times(8.44 + 1.46)} = 10^{3.96} = \num{9120}$.

\textbf{9.} $z = \ang{90} - \ang{82.09} = \ang{7.91}$, $X = \sec\ang{7.91} =
1.0096$；消光为 $0.15\times1.0096 = 0.151$ 星等；测得
$m = -1.46 + 0.151 = -1.31$.

\textbf{10.} $z = \ang{60}$，$X = 2.000$；消光 0.30 星等，
$m = -1.16$。透过的比例为 $10^{-0.4\times0.30} = 0.759$：
损失了 \qty{24}{\percent} 的光。

\textbf{11.} 在升起时 $z = \ang{90}$，$\sec z$ 为无穷大，而
实际穿过弯曲大气的路径长而有限；折射也
使光线弯曲。在地平线附近，消光很大且变化迅速，
而平行平面近似下的修正是错误的。

\textbf{12.} $\num{3.0e4}\times10^{0.4\times(10.0 - 8.44)} = \num{3.0e4}
\times10^{0.624} = \num{3.0e4}\times4.21 = \num{1.26e5}$ 个光电子
每秒。

\textbf{13.} \omterm{def:b1:observing-the-sky:detector}{信噪比}为 1000 时需要 $N = 1000^2 = \num{1.0e6}$
个光电子：$t = \num{1.0e6}/\num{1.26e5} = \qty{7.9}{s}$。星等
误差为 $1.086/1000 = \num{0.0011}$。

\textbf{14.} $\num{3.0e4}\times10^{0.4\times(10.0 + 1.46)} = \num{3.0e4}
\times10^{4.584} = \num{1.15e9}$ 个光电子每秒。

\textbf{15.} $0.20\times\num{1.15e9} = \num{2.30e8}$ 每秒，最亮的
像元中，充满所需时间为 $\num{1.0e5}/\num{2.30e8} =
\qty{4.3e-4}{s}$，不到半毫秒。

\textbf{16.} 天狼星的光翼给出 $\num{1.15e9}\times10^{-4} =
\num{1.15e5}$ 每秒；天狼星~B 给出 $0.20\times\num{1.26e5} =
\num{2.52e4}$ 每秒。天狼星的光晕比天狼星~B 的峰值亮
4.6 倍。

\textbf{17.} 该像元接收 $\num{1.15e5} + \num{2.52e4} = \num{1.40e5}$
个光电子每秒，充满所需时间为 $\num{1.0e5}/\num{1.40e5} = \qty{0.71}{s}$。在
\qty{0.5}{s} 内：信号 $\num{2.52e4}\times0.5 = \num{1.26e4}$，总计
$\num{7.01e4}$，噪声 $\sqrt{\num{7.01e4}} = 265$，因此 $\mathrm{SNR} =
\num{1.26e4}/265 = 48$。天狼星~B 就光子噪声而言是可以探测到的；
真正的困难在于，光晕必须被测定并扣除，精度
需远优于其值的五分之一，而天狼星的核心在
不到一毫秒内即饱和。

\textbf{18.} $m - M = 5\log_{10}(2.637/10) = 5\times(-0.579) = -2.89$.

\textbf{19.} $M = -1.46 + 2.89 = 1.43$.

\textbf{20.} $d = 1/0.13023 = \qty{7.679}{\parsec}$;
$m - M = 5\log_{10}0.7679 = -0.57$，因此 $M = 0.03 + 0.57 = 0.60$。

\textbf{21.} 织女星的\omterm{def:b1:observing-the-sky:absolute}{绝对星等}较小，两者相差的倍数为
$10^{0.4\times(1.43 - 0.60)} = 10^{0.332} = 2.15$。天狼星看上去更亮
只是因为它离我们更近。

\textbf{22.} $m = M_\odot + (m - M) = 4.83 - 2.89 = 1.94$：太阳将会是
一颗相当亮的肉眼可见恒星。

\textbf{23.} $5\log_{10}(d/\qty{10}{\parsec}) = 6.0 - 1.43 = 4.57$，因此
$d = 10\times10^{0.914} = \qty{82.0}{\parsec}$，约 \qty{267}{\lightyear}。

\textbf{24.} $L/L_\odot = 10^{0.4\times(4.83 - 1.43)} = 10^{1.36} = 22.9$:
在 V 波段，天狼星的\omterm{def:b1:observing-the-sky:flux}{光度}相当于\textbf{约 23 个太阳}。
\end{solution}
